\section{仿真验证}

\subsection{仿真设置}

机械臂自由度取 \(n=7\)，惯性参数依据严宇新论文表2-1中的空间机械臂参数整理\cite{yan}。程序实现上将仿真计算与绘图过程分离：仿真脚本先将各算例数据保存为 MAT 文件，绘图脚本再加载该数据统一生成论文插图。由于完整自由漂浮七自由度机械臂动力学需要递推动力学与广义雅可比矩阵构造，程序中采用基于上述参数构造的正定名义关节空间模型作为可运行验证模型。该处理使仿真重点集中在控制律结构、周期延迟反馈实现和扰动下跟踪性能验证上。

仿真步长为 \(T_s=10^{-3}\,\mathrm{s}\)，总时长为 \(10\,\mathrm{s}\)。参考轨迹初始点和终端点分别取
\[
\begin{aligned}
q_r(0)&=[0,\pi/3,0,\pi/4,\pi/4,0,\pi/6]^\top,\\
q_f&=[\pi/18,\pi/6,\pi/5,0,-\pi/4,-\pi/3,\pi/12]^\top .
\end{aligned}
\]
期望轨迹由五次多项式在 \(8\,\mathrm{s}\) 内生成，以保证 \(q_d,\dot q_d,\ddot q_d\) 连续。为避免零初始误差掩盖收敛过程，实际初始关节角设置为
\[
q(0)=q_r(0)+\frac{\pi}{180}[4,-3,3,-4,3,-2,2]^\top,
\qquad
\dot q(0)=0 .
\]

PDF 人工延迟取 \(h=1\,\mathrm{s}\)。激活函数为
\[
R_h(t)=
\begin{cases}
0, & \theta\in[0,h),\\
\sin^4\left(\pi(\theta-h)/h\right), & \theta\in[h,2h),
\end{cases}
\quad \theta=t\bmod 2h .
\]
反馈基准增益取
\[
\begin{aligned}
K_p&=\operatorname{diag}(25,25,25,22,22,20,20),\\
K_d&=\operatorname{diag}(10,10,10,9,9,8,8),
\end{aligned}
\]
并令 \(K_0=[-K_p\ -K_d]\)。PDF 增益 \(K_c(t)\) 按式 \eqref{eq:wc_matrix}--\eqref{eq:kc_pdf} 由离散积分离线计算，所得 Gramian 条件数为 \(5.796\times10^5\)。扰动算例中加入有界时变力矩
\[
d_i(t)=0.15\sin(0.7t+0.3i)+0.05\cos(1.3t+0.2i),
\quad i=1,\ldots,7 .
\]
预设时间扰动观测器参数取 \(T_o=T_c=1\,\mathrm{s}\)、\(\eta=0.3\)、\(\sigma=0.4\)，并采用观测--PDF 两阶段时序。为便于评价，本文给出无延迟计算力矩控制（PD）、PDF 控制以及带预设时间扰动观测器的 PDF 控制（PDF+PTDO）在不同工况下的对照结果。

\subsection{结果与讨论}

图\ref{fig:motion}首先给出了扰动条件下 PDF+PTDO 控制器对应的自由漂浮基座平动、机械臂末端轨迹、初始和最终时刻的机械臂构型以及关节姿态变化。图左侧为三维轨迹与构型示意图，包含基座轨迹、基座初末位置、各关节节点和末端位置；图右侧三行依次给出基座位置、末端位置和关节姿态随时间的变化。其中基座平动由初始总线动量为零时的系统质心守恒关系计算，末端轨迹由随基座平动的名义串联几何链得到。由于当前可运行仿真尚未包含完整角动量递推和基座姿态更新，图中主要展示基座质心平动对机械臂构型和末端轨迹的影响。可以看出，关节运动会诱导基座产生反作用位移，机械臂末端在基座运动和关节运动共同作用下完成空间位移；各关节姿态在五次多项式参考轨迹作用下平滑变化，并在参考轨迹结束后趋于稳定。

\begin{figure}[H]
\centering
\includegraphics[width=\textwidth]{figures/end_effector_joint_position_trajectory.png}
\caption{自由漂浮基座平动、末端轨迹与关节姿态变化}
\label{fig:motion}
\end{figure}

图\ref{fig:joint_summary}进一步给出了扰动条件下 PDF+PTDO 控制器的关节跟踪、误差和控制力矩综合响应。图中第一列按七行分别给出 \(q_1\) 至 \(q_7\) 的位置跟踪曲线，第二列按七行分别给出 \(\dot q_1\) 至 \(\dot q_7\) 的速度跟踪曲线，第三列自上而下分别为关节位置误差、速度误差和控制力矩。由于初始状态与参考初值存在偏差，响应初期出现可见调整过程；随后各关节能够跟随五次多项式参考轨迹，并在参考轨迹结束后保持有界稳定。误差在初始阶段快速下降，在预设时间扰动观测补偿作用下最终收敛到较小范围内；控制力矩未出现指定时间控制中常见的终端奇异增益现象。

\begin{figure}[H]
\centering
\includegraphics[width=\textwidth]{figures/joint_tracking_error_torque_summary.png}
\caption{关节位置与速度跟踪、跟踪误差及控制力矩综合响应}
\label{fig:joint_summary}
\end{figure}

图\ref{fig:observer}给出了等效加速度扰动 \(\Delta_a\) 与观测值 \(\hat\Delta_a\) 的对比，以及观测误差范数 \(\|\tilde\Delta_a\|\)。由于观测器参数选取兼顾峰值抑制和控制输入平滑，扰动估计误差并未在离散数值仿真中达到严格零值，但终端观测误差范数为 \(8.489\times10^{-2}\)，均方根观测误差为 \(3.325\times10^{-2}\)。结合表\ref{tab:sim_metrics}可知，该观测补偿能够显著削弱扰动对误差通道的影响。

\begin{figure}[H]
\centering
\includegraphics[width=\textwidth]{figures/disturbance_observer.png}
\caption{等效加速度扰动真实值、观测值与观测误差}
\label{fig:observer}
\end{figure}

综上，仿真结果验证了所提控制结构的可运行性：在名义模型下，基于 Gramian 构造的光滑周期延迟反馈能够实现高精度跟踪；在有界扰动条件下，预设时间扰动观测补偿能够显著降低终端跟踪误差。需要指出的是，本文使用的名义模型仍是对完整自由漂浮空间机械臂动力学的近似，且观测器参数为抑制峰值进行了保守选取，因此仿真结果主要用于验证控制结构的数值可行性和扰动补偿效果。
